% Options for packages loaded elsewhere
\PassOptionsToPackage{unicode}{hyperref}
\PassOptionsToPackage{hyphens}{url}
\PassOptionsToPackage{dvipsnames,svgnames,x11names}{xcolor}
%
\documentclass[
  11pt,
]{article}
\usepackage{amsmath,amssymb}
\usepackage{iftex}
\ifPDFTeX
  \usepackage[T1]{fontenc}
  \usepackage[utf8]{inputenc}
  \usepackage{textcomp} % provide euro and other symbols
\else % if luatex or xetex
  \usepackage{unicode-math} % this also loads fontspec
  \defaultfontfeatures{Scale=MatchLowercase}
  \defaultfontfeatures[\rmfamily]{Ligatures=TeX,Scale=1}
\fi
\usepackage{lmodern}
\ifPDFTeX\else
  % xetex/luatex font selection
\fi
% Use upquote if available, for straight quotes in verbatim environments
\IfFileExists{upquote.sty}{\usepackage{upquote}}{}
\IfFileExists{microtype.sty}{% use microtype if available
  \usepackage[]{microtype}
  \UseMicrotypeSet[protrusion]{basicmath} % disable protrusion for tt fonts
}{}
\makeatletter
\@ifundefined{KOMAClassName}{% if non-KOMA class
  \IfFileExists{parskip.sty}{%
    \usepackage{parskip}
  }{% else
    \setlength{\parindent}{0pt}
    \setlength{\parskip}{6pt plus 2pt minus 1pt}}
}{% if KOMA class
  \KOMAoptions{parskip=half}}
\makeatother
\usepackage{xcolor}
\usepackage[margin=2.2cm]{geometry}
\setlength{\emergencystretch}{3em} % prevent overfull lines
\providecommand{\tightlist}{%
  \setlength{\itemsep}{0pt}\setlength{\parskip}{0pt}}
\setcounter{secnumdepth}{5}
\usepackage{lmodern}
\usepackage[T1]{fontenc}
\usepackage{amsmath,amssymb}
\usepackage{newunicodechar}
\newunicodechar{}{\ensuremath{\bar\nu}}
\newunicodechar{}{\ensuremath{\hat\nu}}
\newunicodechar{}{\ensuremath{\tilde\sigma}}
\newunicodechar{}{\ensuremath{\bar v}}
\newunicodechar{}{\ensuremath{\hat\sigma}}
\newunicodechar{}{\ensuremath{\bar X}}
\newunicodechar{}{\ensuremath{\hat X}}
\newunicodechar{}{\ensuremath{\tilde W}}
\newunicodechar{}{\ensuremath{\dot\alpha}}
\newunicodechar{}{\ensuremath{\dot\beta}}
\newunicodechar{}{\ensuremath{\hat f}}
\newunicodechar{}{\ensuremath{\bar t}}
\newunicodechar{}{\ensuremath{\hat V}}
\newunicodechar{Ĉ}{\ensuremath{\hat C}}
\newunicodechar{ĉ}{\ensuremath{\hat c}}
\newunicodechar{ġ}{\ensuremath{\dot g}}
\newunicodechar{ᵀ}{\ensuremath{{}^{\mathsf{T}}}}
\newunicodechar{ℝ}{\ensuremath{\mathbb{R}}}
\newunicodechar{−}{\ensuremath{-}}
\newunicodechar{²}{\ensuremath{{}^{2}}}
\newunicodechar{³}{\ensuremath{{}^{3}}}
\newunicodechar{¹}{\ensuremath{{}^{1}}}
\newunicodechar{⁴}{\ensuremath{{}^{4}}}
\newunicodechar{⁸}{\ensuremath{{}^{8}}}
\newunicodechar{⁺}{\ensuremath{{}^{+}}}
\newunicodechar{⁻}{\ensuremath{{}^{-}}}
\newunicodechar{₂}{\ensuremath{{}_{2}}}
\newunicodechar{½}{\ensuremath{\tfrac12}}
\newunicodechar{¼}{\ensuremath{\tfrac14}}
\newunicodechar{·}{\ensuremath{\cdot}}
\newunicodechar{×}{\ensuremath{\times}}
\newunicodechar{÷}{\ensuremath{\div}}
\newunicodechar{±}{\ensuremath{\pm}}
\newunicodechar{σ}{\ensuremath{\sigma}}
\newunicodechar{φ}{\ensuremath{\varphi}}
\newunicodechar{ν}{\ensuremath{\nu}}
\newunicodechar{τ}{\ensuremath{\tau}}
\newunicodechar{π}{\ensuremath{\pi}}
\newunicodechar{ρ}{\ensuremath{\rho}}
\newunicodechar{λ}{\ensuremath{\lambda}}
\newunicodechar{μ}{\ensuremath{\mu}}
\newunicodechar{κ}{\ensuremath{\kappa}}
\newunicodechar{ω}{\ensuremath{\omega}}
\newunicodechar{β}{\ensuremath{\beta}}
\newunicodechar{α}{\ensuremath{\alpha}}
\newunicodechar{γ}{\ensuremath{\gamma}}
\newunicodechar{θ}{\ensuremath{\theta}}
\newunicodechar{η}{\ensuremath{\eta}}
\newunicodechar{χ}{\ensuremath{\chi}}
\newunicodechar{ψ}{\ensuremath{\psi}}
\newunicodechar{ξ}{\ensuremath{\xi}}
\newunicodechar{δ}{\ensuremath{\delta}}
\newunicodechar{ε}{\ensuremath{\varepsilon}}
\newunicodechar{Σ}{\ensuremath{\Sigma}}
\newunicodechar{Δ}{\ensuremath{\Delta}}
\newunicodechar{Λ}{\ensuremath{\Lambda}}
\newunicodechar{Π}{\ensuremath{\Pi}}
\newunicodechar{Ξ}{\ensuremath{\Xi}}
\newunicodechar{Γ}{\ensuremath{\Gamma}}
\newunicodechar{Φ}{\ensuremath{\Phi}}
\newunicodechar{Θ}{\ensuremath{\Theta}}
\newunicodechar{Ψ}{\ensuremath{\Psi}}
\newunicodechar{∫}{\ensuremath{\int}}
\newunicodechar{∂}{\ensuremath{\partial}}
\newunicodechar{√}{\ensuremath{\surd}}
\newunicodechar{∞}{\ensuremath{\infty}}
\newunicodechar{≈}{\ensuremath{\approx}}
\newunicodechar{≥}{\ensuremath{\ge}}
\newunicodechar{≤}{\ensuremath{\le}}
\newunicodechar{≠}{\ensuremath{\ne}}
\newunicodechar{≡}{\ensuremath{\equiv}}
\newunicodechar{→}{\ensuremath{\to}}
\newunicodechar{←}{\ensuremath{\leftarrow}}
\newunicodechar{⇒}{\ensuremath{\Rightarrow}}
\newunicodechar{⇔}{\ensuremath{\Leftrightarrow}}
\newunicodechar{↦}{\ensuremath{\mapsto}}
\newunicodechar{∈}{\ensuremath{\in}}
\newunicodechar{∝}{\ensuremath{\propto}}
\newunicodechar{∏}{\ensuremath{\prod}}
\newunicodechar{⊥}{\ensuremath{\perp}}
\newunicodechar{‖}{\ensuremath{\|}}
\newunicodechar{⟨}{\ensuremath{\langle}}
\newunicodechar{⟩}{\ensuremath{\rangle}}
\newunicodechar{∩}{\ensuremath{\cap}}
\newunicodechar{⊆}{\ensuremath{\subseteq}}
\newunicodechar{…}{\ensuremath{\ldots}}
\newunicodechar{̄}{}
\newunicodechar{̂}{}
\newunicodechar{̃}{}
\newunicodechar{̇}{}
\ifLuaTeX
  \usepackage{selnolig}  % disable illegal ligatures
\fi
\IfFileExists{bookmark.sty}{\usepackage{bookmark}}{\usepackage{hyperref}}
\IfFileExists{xurl.sty}{\usepackage{xurl}}{} % add URL line breaks if available
\urlstyle{same}
\hypersetup{
  pdftitle={Assignment Set 2 --- Defense Sheet},
  pdfauthor={Computational Finance (WI4151), TU Delft --- Prof.~Fang Fang},
  colorlinks=true,
  linkcolor={blue},
  filecolor={Maroon},
  citecolor={Blue},
  urlcolor={Blue},
  pdfcreator={LaTeX via pandoc}}

\title{Assignment Set 2 --- Defense Sheet}
\author{Computational Finance (WI4151), TU Delft --- Prof.~Fang Fang}
\date{}

\begin{document}
\maketitle

{
\hypersetup{linkcolor=blue}
\setcounter{tocdepth}{3}
\tableofcontents
}
\begin{quote}
Scope: pricing a European basket call
\texttt{V(0)\ =\ e\^{}\{-rT\}\ E\^{}Q{[}max(w1\ S1(T)\ +\ w2\ S2(T)\ -\ K,\ 0){]}}
with \(w1 = w2 = 0.5\), \(T = 1\), \(r = 0.02\), \(K = 100\), by three
methods: Q1 Monte Carlo, Q2 moment matching, Q3 the 2D COS method. Two
market setups recur in every question: \textbf{(a)} both underlyings
log-normal (GBM), \(S10=S20=100\), \(sigma1=0.20\), \(sigma2=0.30\),
\(rho=0.30\); \textbf{(b)} \(S1\) log-normal, \(S2\) a CIR / square-root
process \(dS2 = r S2 dt + sigma2 sqrt(S2) dW2\), \(sigma1=0.25\),
\(sigma2=1.5\), \(rho=0.10\) in Q1 / \(rho=0\) in Q2 and Q3.

NOTE: the assignment PDF could not be rendered in this environment (no
\(poppler\)/\(pypdf\)), so the ``task in my own words'' below is
reconstructed from the in-code comments and parameters. \textbf{TODO
(Bruno):} open
\texttt{Computational\_Finance\_Assignment\_2\ (2)\ (1).pdf} and confirm
the exact wording of each sub-question (especially what convergence
behaviour Q1 and Q3 ask you to demonstrate).
\end{quote}

\hypertarget{the-task-in-my-own-words}{%
\section{The task (in my own words)}\label{the-task-in-my-own-words}}

Price a two-asset European \textbf{basket call} under the risk-neutral
measure Q. The basket is \(B_T = w1 S1(T) + w2 S2(T)\); the payoff is
\(max(B_T - K, 0)\), discounted at \(e^{-rT}\).

\begin{itemize}
\tightlist
\item
  \textbf{Q1 --- Monte Carlo.} Simulate the basket and report the price
  with a standard error / 95\% CI.

  \begin{enumerate}
  \def\labelenumi{(\alph{enumi})}
  \tightlist
  \item
    both assets GBM, correlated, exact terminal simulation. (b) \(S1\)
    GBM, \(S2\) a CIR square-root process under the Q-drift \(r S2\);
    demonstrate (i) the weak / time-discretization convergence of the
    Euler scheme and (ii) the statistical \texttt{O(1/sqrt(M))}
    convergence.
  \end{enumerate}
\item
  \textbf{Q2 --- Moment matching.} Compute the raw moments of \(B_T\) in
  closed form, then approximate the basket distribution by matching 2,
  3, 4 moments (log-normal, shifted log-normal, Johnson SU) and price
  the call from each. Benchmark against a high-precision Monte Carlo
  price.
\item
  \textbf{Q3 --- 2D COS.} Price the basket with the two-dimensional COS
  method (Ruijter \& Oosterlee 2012), expanding the joint density of the
  two state variables on \([a1,b1] x [a2,b2]\), exploiting \(rho=0\) so
  the 2D density coefficients factorise into a product of 1D
  coefficients. Show convergence in \(N\).
\end{itemize}

\textbf{TODO (Bruno):} confirm whether the PDF asks for these specific
deliverables or also for, e.g., a comparison table across the three
methods.

\hypertarget{my-approach-and-why}{%
\section{My approach and why}\label{my-approach-and-why}}

\begin{itemize}
\tightlist
\item
  \textbf{Q1:} exact GBM simulation at maturity where possible (no time
  grid needed); Euler--Maruyama with full truncation for the CIR leg
  because exact joint simulation is unavailable when \(rho != 0\).
\item
  \textbf{Q2:} derive \(E[B_T^k]\) analytically from the binomial
  expansion of \((w1 S1 + w2 S2)^k\) and the cross moments of correlated
  log-normals (and the non-central chi-square moments for the CIR leg),
  then fit successively richer distributions to more moments.
\item
  \textbf{Q3:} 1D characteristic functions per marginal, cumulant-rule
  truncation range per marginal, payoff coefficients by 2D trapezoidal
  quadrature, assembled as \(chi1^T V chi2\).
\end{itemize}

\textbf{TODO (Bruno):} for each method, state in one sentence \emph{why}
it is the right tool and what each buys you (MC = unbiased baseline,
moment matching = cheap closed-form approximation, COS = spectral
accuracy). Do not let the examiner hear a generic answer.

\hypertarget{key-code-sections-file-what-it-does-why-i-wrote-it-that-way}{%
\section{Key code sections (file + what it does + why I wrote it that
way)}\label{key-code-sections-file-what-it-does-why-i-wrote-it-that-way}}

\hypertarget{q1-basket-monte-carlo-q1_basket_mc-1.py}{%
\subsection{\texorpdfstring{Q1 --- basket Monte Carlo
(\texttt{Q1\_basket\_MC\ (1).py})}{Q1 --- basket Monte Carlo (Q1\_basket\_MC (1).py)}}\label{q1-basket-monte-carlo-q1_basket_mc-1.py}}

\begin{itemize}
\tightlist
\item
  \textbf{\texttt{simulate\_basket\_lognormal(M)} --- lines 17--42.}
  Q1(a). Draws two standard normals (\(Z1\), \(Zc\), lines 27--28) and
  correlates them by Cholesky: \(Z2 = rho*Z1 + sqrt(1-rho^2)*Zc\) (line
  29). Simulates terminal prices with the \textbf{exact GBM solution} at
  \(T\) (lines 32--33), forms the basket and discounted payoff (lines
  36--37), and returns the MC mean and standard error
  \texttt{std(ddof=1)/sqrt(M)} (lines 40--41).
\item
  \textbf{Convergence loop --- lines 48--63.} Prices over
  \(M = 2000 … 1e6\) (line 48), then estimates the convergence rate as
  the OLS slope of \(log(SE)\) vs \(log(M)\) (line 62), expected
  \(-0.5\).
\item
  \textbf{\texttt{simulate\_basket\_cir(M,\ n\_steps)} --- lines
  93--125.} Q1(b). Time-steps with \texttt{h\ =\ T/n\_steps} (line 107).
  \(S1\) exact GBM step (line 118); \(S2\) \textbf{Euler--Maruyama} step
  with \textbf{full truncation} \(S2p = max(S2,0)\) inside drift and
  diffusion (lines 120--121), so the square root is always defined.
\item
  \textbf{Weak-error test --- lines 134--185.} Reuses the \textbf{same
  fine Brownian increments} across all coarse grids
  (\texttt{basket\_on\_grid}, lines 144--159; increments summed into
  blocks at lines 151--152) to isolate the discretization bias from MC
  noise. Fits the weak order as the slope of \(log|bias|\) vs \(log h\)
  on the coarse grids (line 184), expected \(1.0\).
\item
  \textbf{Statistical-error test --- lines 198--210.} Fixes the grid,
  varies \(M\), fits the SE slope (line 209).
\end{itemize}

\textbf{TODO (Bruno):} justify (i) full truncation vs reflection / other
CIR fixes; (ii) why you reuse the same increments across grids for the
weak-error test; (iii) the Q-drift argument at lines 86--91 (physical
CIR drift \(kappa(theta - S2)\) replaced by \(r S2\) because the
discounted price is a Q-martingale).

\hypertarget{q2-moment-matching-q2_moment_matching-1.py}{%
\subsection{\texorpdfstring{Q2 --- moment matching
(\texttt{Q2\_moment\_matching\ (1).py})}{Q2 --- moment matching (Q2\_moment\_matching (1).py)}}\label{q2-moment-matching-q2_moment_matching-1.py}}

\begin{itemize}
\tightlist
\item
  \textbf{\(lognormal_cross(j, m, ...)\) --- lines 20--29.} Cross raw
  moment \texttt{E{[}S1\^{}j\ S2\^{}m{]}} for two correlated
  log-normals, from \(E[exp(N)] = exp(mean + 0.5 var)\) with the
  combined log-return mean/variance/cov at lines 26--29.
\item
  \textbf{\texttt{basket\_moments\_LN(...)} --- lines 31--37.} Raw
  basket moments \(E[B_T^k]\) via the binomial expansion
  \texttt{sum\_j\ C(k,j)\ w1\^{}j\ w2\^{}(k-j)\ E{[}S1\^{}j\ S2\^{}(k-j){]}}.
\item
  \textbf{\texttt{cir\_S2\_rawmoments(...)} --- lines 39--53.} Raw
  moments of \(S2(T) = c·Z\), \(Z ~ chi'^2_0(lambda)\) (zero-dof
  non-central chi-square), using cumulants
  \texttt{kappa\_n\ =\ 2\^{}(n-1)\ n!\ lambda} (line 49) and the
  raw-from-cumulant recursion (lines 50--52). \(c\) and \texttt{lambda}
  from the technical hint at lines 46--48.
\item
  \textbf{\texttt{basket\_moments\_CIR(...)} --- lines 55--64.} Combines
  the LN-\(S1\) and CIR-\(S2\) marginal moments under independence
  (\(rho=0\)), so joint moments factorise (line 63).
\item
  \textbf{\texttt{central\_from\_raw(mom)} --- lines 66--75.} Converts
  raw moments to mean, variance, skewness, kurtosis.
\item
  \textbf{\(price_2moments\) --- lines 78--89.} Fits a log-normal
  matching \(m1\), \(m2\); Black's formula with \(F = m1\), total
  variance \texttt{s\^{}2\ =\ log(m2/m1\^{}2)}.
\item
  \textbf{\(price_3moments\) --- lines 91--107.} Shifted log-normal
  \(B = gamma + L\); inverts the log-normal skewness
  \((w+2)sqrt(w-1) = g1\) for \(w\) by \(brentq\) (line 99), recovers
  shift and shifted strike.
\item
  \textbf{\(price_4moments\) --- lines 109--139.} Johnson SU; solves the
  2×2 system for \((delta, gamma)\) from skewness and kurtosis by
  \(least_squares\) with Gauss--Hermite moments of \(sinh\)
  (\(U_moments\), lines 120--125; residual lines 127--132; solve line
  134), then prices by Gauss--Hermite quadrature (line 139). GH
  nodes/weights set at lines 16--17.
\item
  \textbf{MC benchmarks --- \texttt{mc\_benchmark\_LN} lines 144--159,
  \texttt{mc\_benchmark\_CIR} lines 161--180.} Exact simulation with
  antithetic variates; the CIR benchmark samples the zero-dof
  non-central chi-square via a Poisson mixture of Gammas (lines
  171--173).
\end{itemize}

\textbf{TODO (Bruno):} be ready to (i) derive the cross-moment formula
at line 29 by hand; (ii) derive the non-central chi-square cumulants
\texttt{kappa\_n\ =\ 2\^{}(n-1)\ n!\ lambda}; (iii) explain \emph{why}
more matched moments need not monotonically reduce the pricing error.

\hypertarget{q3-basket-cos-q3_basket_cos-1.py}{%
\subsection{\texorpdfstring{Q3 --- basket COS
(\texttt{Q3\_basket\_COS\ (1).py})}{Q3 --- basket COS (Q3\_basket\_COS (1).py)}}\label{q3-basket-cos-q3_basket_cos-1.py}}

\begin{itemize}
\tightlist
\item
  \textbf{\(chi_coeffs(phi, a, b, N)\) --- lines 22--32.} 1D density
  cosine coefficients
  \texttt{chi\_k\ =\ Re\{\ phi(omega\_k)\ exp(-i\ omega\_k\ a)\ \}},
  \texttt{omega\_k\ =\ k\ pi/(b-a)} (lines 29--30), with the \(k=0\)
  term halved (line 31) to realise the primed summation.
\item
  \textbf{\texttt{payoff\_coeffs\_2d(...)} --- lines 34--54.} Payoff
  cosine coefficients \(V_{k1,k2}\): builds the payoff matrix
  \(P = max(w1 g1(y1) + w2 g2(y2) - K, 0)\) on a dense grid (line 49)
  and contracts it with cosine/trapezoidal-weight matrices \(C1, C2\)
  (lines 52--53) as \(C1^T P C2\) (line 54). Trapezoidal weights at
  lines 45--46; normalisation \texttt{(2/(b1-a1))(2/(b2-a2))} at line
  54.
\item
  \textbf{\(price_2dCOS(...)\) --- lines 56--63.} Assembles
  \texttt{e\^{}\{-rT\}\ chi1\^{}T\ V\ chi2} (line 63).
\item
  \textbf{\(cumulant_range(k1, k2, k4, L=12)\) --- lines 65--71.}
  \textbf{The truncation range.} Returns \([k1 - width, k1 + width]\)
  with \(width = L*sqrt(|k2| + sqrt(|k4|))\) --- \textbf{lines 70--71}.
  This is the COS-paper cumulant rule; \(L=12\) is the default.
\item
  \textbf{Q3(a) characteristic functions --- lines 79--80.} Normal
  log-price CFs
  \texttt{phi(om)\ =\ exp(i\ om\ mu\ -\ 0.5\ sigma\^{}2\ T\ om\^{}2)}
  for each marginal; means at lines 76--77.
\item
  \textbf{Q3(a) ranges --- lines 82--83.}
  \(cumulant_range(mu, sigma^2 T, 0)\) per marginal (a normal has
  \(k1 = mean\), \(k2 = variance\), \(k4 = 0\)).
\item
  \textbf{Q3(b) \(S2\) characteristic function --- \(phi2_b\), lines
  117--119.} CF of \(S2(T) = c·Z\), \(Z ~ chi'^2_0(lambda)\):
  \texttt{exp(lambda\ i\ z\ /\ (1\ -\ 2i\ z))} with \(z = om·c\) (lines
  118--119); \(c_b\), \(lam_b\) at lines 114--116.
\item
  \textbf{Q3(b) cumulants of \(S2\) --- lines 122--124.}
  \texttt{k1\ =\ c·lambda}, \texttt{k2\ =\ c\^{}2·4·lambda},
  \texttt{k4\ =\ c\^{}4·192·lambda} (from
  \texttt{kappa\_n\ =\ 2\^{}(n-1)\ n!\ lambda} scaled by \(c^n\)).
\item
  \textbf{Q3(b) ranges --- lines 125--127.} \(S1\) range
  \texttt{cumulant\_range(mu1\_b,\ sigma1\^{}2\ T,\ 0)} (line 125);
  \(S2\) range \texttt{cumulant\_range(k1\_S2,\ k2\_S2,\ k4\_S2,\ L=12)}
  (line 126), then \textbf{clamped to be non-negative}
  \texttt{a2\_b\ =\ max(a2\_b,\ 0.0)} (line 127) because \(S2 >= 0\).
\item
  \textbf{Payoff maps --- lines 84--85, 128--129.} \(g = exp\) for
  log-price state variables; \(g2_b = identity\) in Q3(b) because the
  \(S2\) state variable is the \emph{level}, not the log.
\item
  \textbf{Convergence loops --- lines 88--104 (a), 131--145 (b).}
  Compare COS prices over \(N = 8 … 64\) against a converged \(N = 256\)
  reference.
\end{itemize}

\hypertarget{design-decisions-i-must-justify}{%
\section{Design decisions I must
justify}\label{design-decisions-i-must-justify}}

\begin{itemize}
\tightlist
\item
  \textbf{Cholesky correlation} (\(Z2 = rho Z1 + sqrt(1-rho^2) Zc\)) ---
  Q1 lines 29, 116; Q2 line 152. \textbf{TODO (Bruno):} why this is
  exactly the 2×2 Cholesky factor.
\item
  \textbf{Full truncation} of the CIR scheme --- Q1 lines 120--121.
  \textbf{TODO (Bruno):} why truncation (not reflection) and what bias
  it introduces.
\item
  \textbf{Reusing fine Brownian increments across grids} for the
  weak-error test --- Q1 lines 144--169. \textbf{TODO (Bruno):} why this
  cleanly isolates discretization bias from MC noise.
\item
  \textbf{Antithetic variates} in the Q2 benchmarks --- Q2 lines
  154--158, 176--179. \textbf{TODO (Bruno):} why antithetics reduce
  variance here and when they fail.
\item
  \textbf{Choice of approximating families} (LN / shifted-LN / Johnson
  SU) --- Q2 lines 78--139. \textbf{TODO (Bruno):} why each family is
  the natural one for 2 / 3 / 4 matched moments.
\item
  \textbf{\(rho = 0\) factorisation} of the 2D density coefficients ---
  Q3 lines 13--20, 60--63. \textbf{TODO (Bruno):} show why independence
  makes \texttt{c\_\{k1,k2\}\ =\ chi\_\{k1\}\ chi\_\{k2\}} and what
  changes if \(rho != 0\).
\item
  \textbf{\(L = 12\) in the cumulant rule} --- Q3 line 65, used at lines
  82--83, 125--126. \textbf{TODO (Bruno):} justify \(L = 12\) (the paper
  often uses 8--10); relate to the bug below.
\item
  \textbf{Trapezoidal quadrature for \(V_{k1,k2}\)} instead of a closed
  form --- Q3 lines 34--54. \textbf{TODO (Bruno):} why no analytic
  \(V_k\) here (basket payoff couples the two state variables), and how
  \(Ngrid\) was chosen.
\end{itemize}

\hypertarget{results-and-what-they-mean}{%
\section{Results and what they mean}\label{results-and-what-they-mean}}

\textbf{TODO (Bruno):} run the three scripts and fill in the actual
numbers: - Q1(a) price + 95\% CI, and the fitted SE slope (should be ≈
\(-0.5\), line 63). - Q1(b) price + CI, the fitted weak order (≈
\(1.0\), line 184), the fitted SE slope (line 209). - Q2(a) and Q2(b):
basket mean/std/skew/kurtosis, MC benchmark, and the 2/3/4-moment prices
with absolute errors (printed at lines 196--198, 214--216) --- explain
why error generally falls with more moments and any exception. -
Q3(a)/(b): converged COS price and the error/ratio table (lines 98--104,
140--145) --- explain whether convergence looks spectral (exponential)
and how the ratios behave.

\hypertarget{the-hardest-question-fang-could-ask-here-and-my-answer}{%
\section{The hardest question Fang could ask here --- and my
answer}\label{the-hardest-question-fang-could-ask-here-and-my-answer}}

\textbf{Most likely deep-dive (highest probability): the COS truncation
range \([a,b]\).} Fang co-invented COS; the cumulant truncation rule is
her result, and this assignment had a real bug there. She will go
straight to \textbf{\texttt{Q3\_basket\_COS\ (1).py:65-71}}
(\(cumulant_range\)) and its call sites \textbf{\texttt{:82-83} (Q3a)}
and \textbf{\texttt{:125-127} (Q3b)}.

What the code factually does: \(cumulant_range\) returns
\([k1 - L*sqrt(|k2| + sqrt(|k4|)), k1 + L*sqrt(|k2| + sqrt(|k4|))]\)
(\texttt{Q3\_basket\_COS\ (1).py:70-71}), with \(L = 12\). For a normal
marginal \(k4 = 0\), so the half-width is \(L*sqrt(variance)\)
(\texttt{:82-83}). For the CIR leg, the half-width uses the non-central
chi-square cumulants \texttt{k2\ =\ 4\ c\^{}2\ lambda},
\texttt{k4\ =\ 192\ c\^{}4\ lambda} (\texttt{:123-124}), and the lower
end is clamped to \(0\) (\texttt{:127}).

\textbf{The bug (special focus): a too-small \(c2\) made \([a,b]\) too
narrow at short maturities.} The half-width scales like \(sqrt(k2)\),
and \(k2\) is the variance term --- for a normal marginal
\(k2 = sigma^2 T\) (\texttt{:82-83}, \texttt{:125}), which
\textbf{shrinks with \(T\)}. At short maturity \(c2\) is small, so
\(width = L*sqrt(c2 + sqrt(c4))\) (\texttt{Q3\_basket\_COS\ (1).py:70})
collapses and \([a,b]\) becomes too narrow to contain the support of the
density / payoff, truncating mass and corrupting the price. The fix
lives at the range construction --- the \(cumulant_range\) formula at
\textbf{\texttt{:70}} and the \(L\) value at \textbf{\texttt{:65}} (and
the non-negativity clamp at \texttt{:127}).

\begin{quote}
\textbf{TODO (Bruno): the \emph{reasoning} of the fix --- fill this in
yourself.} Specifically: 1. What symptom did you see (wrong/unstable
price, or error not converging) and at which \(T\)? 2. What exactly did
you change --- raise \(L\), add the \(sqrt(k4)\) term, widen the floor,
or change how \(c2\) is fed in for short \(T\)? Point to the precise
line you edited (\texttt{:65} \(L=12\), \texttt{:70} the width formula,
or \texttt{:127} the clamp). 3. Why \(L*sqrt(c2 + sqrt(c4))\) (and not
just \(L*sqrt(c2)\)) is the right safeguard when \(c2\) is tiny but
\(c4 > 0\), and why a fixed wide interval is \emph{not} a free lunch
(more terms \(N\) needed for the same accuracy, since COS error depends
on both truncation and series length).
\end{quote}

This is the single question most worth over-preparing.

\hypertarget{show-me-in-your-code-where-you-do-x-anticipated-spots}{%
\section{``Show me in your code where you do X'' --- anticipated
spots}\label{show-me-in-your-code-where-you-do-x-anticipated-spots}}

\begin{itemize}
\tightlist
\item
  \textbf{``\ldots where you compute the truncation range \([a,b]\).''}
  \texttt{Q3\_basket\_COS\ (1).py:65-71} (\(cumulant_range\)), called at
  \texttt{:82-83} (Q3a) and \texttt{:125-127} (Q3b, with the
  \(max(a2_b, 0)\) clamp at \texttt{:127}).
\item
  \textbf{``\ldots where a small \(c2\) makes \([a,b]\) too narrow, and
  how you fixed it.''} \texttt{Q3\_basket\_COS\ (1).py:70}
  (\(width = L*sqrt(abs(k2) + sqrt(abs(k4)))\)) and \texttt{:65}
  (\(L=12\)). \textbf{TODO (Bruno):} name the line you edited.
\item
  \textbf{``\ldots where you define the characteristic function.''}
  Q3(a) \texttt{:79-80}; Q3(b) CIR leg \(phi2_b\) \texttt{:117-119}.
\item
  \textbf{``\ldots where you build the cosine density coefficients
  \(chi_k\).''} \(chi_coeffs\) \texttt{Q3\_basket\_COS\ (1).py:22-32}
  (note the \(chi[0] *= 0.5\) primed-sum term at \texttt{:31}).
\item
  \textbf{``\ldots where you build the payoff coefficients
  \(V_{k1,k2}\).''} \texttt{payoff\_coeffs\_2d} \texttt{:34-54}
  (assembly \texttt{:54}).
\item
  \textbf{``\ldots where the 2D price is assembled.''} \(price_2dCOS\)
  \texttt{:56-63} (\(chi1 @ V @ chi2\) at \texttt{:63}).
\item
  \textbf{``\ldots where the cumulants of the CIR leg come from.''}
  \texttt{Q3\_basket\_COS\ (1).py:122-124}; same moments in
  \texttt{Q2\_moment\_matching\ (1).py:39-53}.
\item
  \textbf{``\ldots where you correlate the Brownian motions.''}
  \texttt{Q1\_basket\_MC\ (1).py:29} and \texttt{:116}.
\item
  \textbf{``\ldots where you do the Euler step for \(S2\).''}
  \texttt{Q1\_basket\_MC\ (1).py:120-121} (full truncation).
\item
  \textbf{``\ldots where you isolate the discretization bias.''}
  \texttt{Q1\_basket\_MC\ (1).py:144-169}.
\item
  \textbf{``\ldots where you match 2 / 3 / 4 moments.''}
  \texttt{price\_2moments\ :78-89}, \texttt{price\_3moments\ :91-107},
  \texttt{price\_4moments\ :109-139} in
  \texttt{Q2\_moment\_matching\ (1).py}.
\item
  \textbf{``\ldots where the MC benchmark uses antithetics.''}
  \texttt{Q2\_moment\_matching\ (1).py:154-158} and \texttt{:176-179}.
\end{itemize}

\begin{center}\rule{0.5\linewidth}{0.5pt}\end{center}

\hypertarget{corrector-feedback-from-the-graded-pdf}{%
\section{Corrector feedback (from the graded
PDF)}\label{corrector-feedback-from-the-graded-pdf}}

\begin{quote}
The returned PDF in this folder
(\texttt{Computational\_Finance\_Assignment\_2\ (2)\ (1).pdf}) contains
\textbf{no grader annotations} --- no embedded scores or comments were
found, so there is nothing to reconcile here. The COS truncation-range
discussion above stands on its own. If a separately-graded copy exists,
share it and I'll fold the comments in the same way as for the other
three sets.
\end{quote}

\end{document}
